\documentclass[11pt]{article}
\usepackage{amsmath,amssymb,amsthm,hyperref}
\newtheorem{lemma}{Lemma}
\title{The length-three branch of the ratio-25 problem}
\author{Research attempt prepared in the accompanying Codex conversation}
\date{6 November 2026}
\begin{document}
\maketitle

\noindent
Status: this treats only $k=3$ of the
\href{https://conjectures.io/problems/erdos686-erdos-686-variants-twenty-five}{listed ratio-25 target}.
It does not settle the target, whose block length is unbounded. No novelty
claim is made. The finite classification below depends on a Magma Thue
computation; its completeness has not been independently certified here.

\section*{Reduction to sixty Thue equations}
At length three the required product identity becomes, with
$x=m+2$ and $y=n+2$,
\[
 x^3-x=25(y^3-y),\qquad x,y\ge2.
\]
The disjointness condition is not needed for the exclusion below.

\begin{lemma}
Any positive integer solution with $x,y\ge2$ gives coprime positive
integers $a,b$, a positive divisor $h$ of $15600$, and an integer $g\ge1$
satisfying
\[
 a^3-25b^3=-h,\qquad g^2h=25b-a,\qquad (x,y)=(ga,gb).
\]
\end{lemma}
\begin{proof}
Put $g=\gcd(x,y)$, $x=ga$, $y=gb$. Dividing the cubic identity by $g$
gives
\[
 g^2(a^3-25b^3)=a-25b.
\]
For $t\ge1$, the function $t^3-t$ is strictly increasing, and
$(25y)^3-25y>25(y^3-y)$. Hence $x<25y$ and $a<25b$.
Thus $h=25b^3-a^3$ is positive and $g^2h=25b-a$.
In particular $a\equiv25b\pmod h$. Substitution into
$a^3-25b^3\equiv0\pmod h$ gives
$h\mid25(25^2-1)b^3=15600b^3$.
Since $\gcd(a,b)=1$, also $\gcd(h,b)=1$, so $h\mid15600$.
\end{proof}

\section*{Observed complete-solver output and postfilter}
The script \texttt{erdos686-thue-check.m} was evaluated in the official
\href{https://magma.maths.usyd.edu.au/calc/}{Magma calculator}, version
V2.29-10, with seed 2399912018. The observed output is transcribed in
\texttt{erdos686-thue-results.json}. It computes
\texttt{Solutions(Thue(t\^{}3-25),-h)} for every positive $h\mid15600$.
Magma's \href{https://magma.maths.usyd.edu.au/magma/handbook/text/418}{documentation}
specifies that this returns all integer solutions; no GRH class-group
setting was requested. The
\href{https://magma.maths.usyd.edu.au/magma/handbook/text/416}{class-group documentation}
states that its default proof bound is unconditional.

Among the reported solutions, the only coprime pairs with $a,b>0$ are
\[
\begin{array}{c|c|c|c}
 a&b&h&(25b-a)/h\\\hline
 1&1&24&1\\
 5&2&75&3/5\\
 35&12&325&53/65\\
 25&9&2600&1/13\\
 15&7&5200&2/65
\end{array}
\]
The last four ratios are below one and cannot equal $g^2$ for $g\ge1$.
The first gives $g=1$ and $x=y=1$, outside the required range.
Thus the computed classification excludes length three for every pair of
natural offsets, even without disjointness.

The local script \texttt{erdos686-thue-verify.mjs} checks all sixty right-hand
sides, every transcribed solution, and the exact postfilter. It checks
neither the completeness of the Thue output nor the internal correctness
of Magma's algorithm. A separate bounded elliptic-group search also found
only the trivial cubic pairs; it is not used for completeness.

\section*{Elliptic model used for exploration}
The rational substitutions
\[
 X=625+\frac{15600y}{x-25y},\qquad Y=\frac{15600}{x-25y}
\]
give $Y^2=X^3-1875X+391250$, with inverse
\[
 x=\frac{25(X-1)}Y,\qquad y=\frac{X-625}Y.
\]
Magma reported rank bounds $[2,2]$ and generators $(1,624),(25,600)$,
followed by two true verification flags. No rank computation is needed
for the sixty-equation reduction above. The local exact rational code is
\texttt{erdos686-elliptic-search.mjs}; its completed coefficient box is
$[-12,12]^2$. A larger box was stopped after the Thue computation made
further subgroup searching unnecessary for this branch.
\end{document}
